% AI-era Theory Note — one-page worked template.
% The design records a change in judgment, not just the final proof.
% Compile with XeLaTeX or LuaLaTeX.
\documentclass[10pt,fontset=none]{ctexart}

\usepackage{iftex}
\ifPDFTeX\errmessage{Please compile with XeLaTeX or LuaLaTeX}\fi
\usepackage{fontspec}
\usepackage{geometry}
\usepackage{xcolor}
\usepackage{microtype}
\usepackage{hyperref}
\usepackage{fancyhdr}
\usepackage[most]{tcolorbox}
\usepackage{amsmath,amssymb,mathtools,bm}
\usepackage{tabularx,array}
\usepackage{enumitem}

\geometry{
  paperwidth=7.5in,paperheight=10in,
  left=0.58in,right=0.58in,top=0.34in,bottom=0.40in,
  headsep=0in,footskip=0.18in
}

\IfFontExistsTF{Songti SC}
  {\setCJKmainfont{Songti SC}[AutoFakeBold=2.2]}
  {\IfFontExistsTF{Source Han Serif SC}
    {\setCJKmainfont{Source Han Serif SC}[AutoFakeBold=2.2]}
    {\IfFontExistsTF{Noto Serif CJK SC}
      {\setCJKmainfont{Noto Serif CJK SC}[AutoFakeBold=2.2]}
      {\setCJKmainfont{FandolSong-Regular.otf}[BoldFont=FandolSong-Bold.otf]}}}
\IfFontExistsTF{Times New Roman}
  {\setmainfont{Times New Roman}}
  {\setmainfont{TeX Gyre Termes}}
\IfFontExistsTF{Helvetica Neue}
  {\setsansfont{Helvetica Neue}}
  {\setsansfont{TeX Gyre Heros}}
\IfFontExistsTF{Hiragino Sans GB}
  {\setCJKsansfont{Hiragino Sans GB}[AutoFakeBold=2.2]}
  {\IfFontExistsTF{Noto Sans CJK SC}
    {\setCJKsansfont{Noto Sans CJK SC}[AutoFakeBold=2.2]}
    {\setCJKsansfont{FandolHei-Regular.otf}[BoldFont=FandolHei-Bold.otf]}}

\definecolor{navy}{HTML}{173B78}
\definecolor{blue}{HTML}{3157C8}
\definecolor{blueSoft}{HTML}{F0F4FF}
\definecolor{green}{HTML}{177C68}
\definecolor{greenSoft}{HTML}{EEF8F4}
\definecolor{amber}{HTML}{C97820}
\definecolor{amberSoft}{HTML}{FFF7ED}
\definecolor{ink}{HTML}{171B22}
\definecolor{muted}{HTML}{66717F}
\definecolor{line}{HTML}{D7DEE8}
\definecolor{paper}{HTML}{FBFAF7}

\pagecolor{paper}
\color{ink}
\linespread{1.03}
\setlength{\parindent}{0pt}
\setlength{\parskip}{0pt}
\setlength{\abovedisplayskip}{4pt}
\setlength{\belowdisplayskip}{4pt}
\setlength{\abovedisplayshortskip}{3pt}
\setlength{\belowdisplayshortskip}{3pt}
\hypersetup{colorlinks=true,linkcolor=blue,urlcolor=blue,
  pdfauthor={MIStatlE},pdftitle={AI-era Theory Note: Gradient Descent Contraction}}

\fancyhf{}
\renewcommand{\headrulewidth}{0pt}
\fancyfoot[L]{\small\rmfamily\color{muted}TN--01 · 可复现的判断，才是积累}
\fancyfoot[R]{\small\sffamily\color{blue}@MIStatlE · 1}
\pagestyle{fancy}

\tcbset{
  principle/.style={enhanced,colback=blueSoft,colframe=blue!28,boxrule=0.5pt,
    arc=1pt,left=9pt,right=9pt,top=5pt,bottom=5pt},
  theorem/.style={enhanced,colback=white,colframe=navy!42,boxrule=0.65pt,
    borderline west={3pt}{0pt}{navy},arc=1pt,
    left=9pt,right=9pt,top=6pt,bottom=6pt},
  audit/.style={enhanced,frame hidden,borderline west={3pt}{0pt}{#1},
    colback=#1!6,arc=1pt,left=8pt,right=8pt,top=5pt,bottom=5pt},
  ledger/.style={enhanced,colback=white,colframe=line,boxrule=0.55pt,
    arc=1pt,left=8pt,right=8pt,top=5pt,bottom=5pt}
}

\newcommand{\Kicker}[2]{%
  {\rmfamily\bfseries\scriptsize\color{#1}#2}\par\vspace{0.18em}}

\newcommand{\Status}[2]{%
  \tcbox[on line,boxsep=1pt,left=4pt,right=4pt,top=1pt,bottom=1pt,
    colback=#1!10,colframe=#1!45,boxrule=0.45pt,arc=2pt,
    fontupper=\sffamily\bfseries\scriptsize\color{#1},before upper=\strut]{#2}}

\newcommand{\SectionRule}[2]{%
  \vspace{0.28em}
  \noindent{\rmfamily\bfseries\large\color{#1}#2}\hspace{0.7em}%
  {\color{line}\leaders\hrule height 0.45pt\hfill\kern0pt}\par
  \vspace{0.18em}}

\newcommand{\good}{\textcolor{green}{\ensuremath{\checkmark}}}
\newcommand{\pending}{\textcolor{amber}{\ensuremath{\circ}}}

\begin{document}

\begin{minipage}[t]{0.71\linewidth}
  \Kicker{blue}{THEORY NOTE · JUDGMENT TRACE}
  {\bfseries\fontsize{24pt}{27pt}\selectfont\color{navy}
    梯度下降的收缩率\par}
  \vspace{0.12em}
  {\large\color{muted}不是保存答案，而是保存判断如何被证据改变}
\end{minipage}\hfill
\begin{minipage}[t]{0.25\linewidth}
  \raggedleft
  {\sffamily\scriptsize\color{muted}NOTE 01}\par
  \vspace{0.18em}{\color{line}\rule{0.78\linewidth}{0.5pt}}\par
  \vspace{0.18em}{\sffamily\small\color{green}\textbf{JUDGMENT}\quad REVISED}\par
  \vspace{0.15em}{\sffamily\small\color{muted}NEXT\quad REBUILD ALONE}
\end{minipage}

\vspace{0.28em}
\begin{tcolorbox}[principle]
  \Kicker{blue}{NOTE PRINCIPLE}
  理论笔记不以结论为最小单位，而记录一次\textbf{可审计的判断变化}：
  原判断 $\rightarrow$ 反证信号 $\rightarrow$ AI 压力测试 $\rightarrow$ 人的决定 $\rightarrow$ 独立重建。
\end{tcolorbox}

\vspace{0.08em}
\begin{tcolorbox}[theorem]
  \Kicker{navy}{目标命题 \,/\, TARGET}
  设 $f:\mathbb R^d\to\mathbb R$ 为 $\mu$-强凸且 $L$-光滑，$x^\star$ 为极小点。取
  $\eta=2/(L+\mu)$，目标是保留最优条件数依赖：
  \[
    \|x_{t+1}-x^\star\|
    \le \frac{L-\mu}{L+\mu}\,\|x_t-x^\star\|.
  \]
\end{tcolorbox}

\SectionRule{navy}{1\quad 起点：我原本相信什么}

令 $d_t=x_t-x^\star$，$g_t=\nabla f(x_t)-\nabla f(x^\star)$。第一判断：分别调用强单调性与 Lipschitz 性应该足够。
\[
\|d_t-\eta g_t\|^2
\le \bigl(1-2\eta\mu+\eta^2L^2\bigr)\|d_t\|^2.
\]
这条推导\textbf{正确，却不够好}：最优步长只有 $\eta=\mu/L^2$，把应有的 $\kappa=L/\mu$ 依赖损失成近似 $\kappa^2$。

\vspace{0.22em}
\noindent
\begin{minipage}[t]{0.475\linewidth}
  \begin{tcolorbox}[audit=amber]
    \Kicker{amber}{反证信号 \,/\, EVIDENCE}
    $\kappa^2$ 依赖直接否定了第一判断：分别正确的界，组合后未必保留问题的几何结构。
  \end{tcolorbox}
\end{minipage}\hfill
\begin{minipage}[t]{0.505\linewidth}
  \begin{tcolorbox}[audit=blue]
    \Kicker{blue}{AI 的位置 \,/\, PRESSURE TEST}
    AI 不继续计算，只追问：
    \begin{itemize}[leftmargin=1.25em,itemsep=0.12em,topsep=0.2em]
      \item 两个不等式能否同时取紧？
      \item 是否存在同时保留 $\langle g,d\rangle$ 与 $\|g\|^2$ 的混合量？
    \end{itemize}
  \end{tcolorbox}
\end{minipage}

\vspace{0.08em}
\begin{tcolorbox}[audit=green]
  \Kicker{green}{我的决定 \,/\, HUMAN JUDGMENT}
  放弃“分别取界”的表示，改为保留耦合。采用光滑强凸函数的插值不等式
  \[
  \langle g,d\rangle\ge
  \frac{\mu L}{\mu+L}\|d\|^2+
  \frac{1}{\mu+L}\|g\|^2.
  \]
\end{tcolorbox}

\SectionRule{green}{2\quad 重建：保留耦合的证明主线}

\begin{enumerate}[leftmargin=1.7em,label=\textbf{\arabic*.},itemsep=0.25em,topsep=0pt]
  \item 代入恒等式 $\|d_t-\eta g_t\|^2=\|d_t\|^2-2\eta\langle g_t,d_t\rangle+\eta^2\|g_t\|^2$。
  \item 取 $\eta=2/(L+\mu)$，$\|g_t\|^2$ 的系数恰好抵消。
  \item 剩余系数为
  \[
    1-\frac{4\mu L}{(L+\mu)^2}
    =\left(\frac{L-\mu}{L+\mu}\right)^2,
  \]
  开方即得目标收缩率。
\end{enumerate}

\SectionRule{navy}{3\quad 能力增量：下次我会做出什么不同的第一步}

\begin{tcolorbox}[ledger]
  \small
  \begin{minipage}[t]{0.30\linewidth}
    \Kicker{green}{判断已改变}
    以后遇到“局部正确、整体过松”的证明，先检查表示是否拆得太早。
  \end{minipage}\hfill
  \begin{minipage}[t]{0.28\linewidth}
    \Kicker{navy}{证据已核验}
    \good\ 代数消项与最终系数\par
    \vspace{0.16em}\good\ 边界 $\mu=L$ 一致
  \end{minipage}\hfill
  \begin{minipage}[t]{0.36\linewidth}
    \Kicker{amber}{仍欠自己的证明}
    \pending\ 不看 AI 对话，从 $f-\frac\mu2\|\cdot\|^2$ 的凸光滑性重建插值引理。
  \end{minipage}
\end{tcolorbox}

\end{document}
